% !TEX program = xelatex
\documentclass{beamer}

\usepackage[UTF8]{ctex}
\usepackage{graphicx}
\usepackage{booktabs}
\usepackage{amsmath}
\usepackage{array}

\usetheme{default}
\usecolortheme{default}
\setbeamertemplate{navigation symbols}{}
\setbeamertemplate{footline}[frame number]

\newcommand{\countyfig}[2][]{%
  \includegraphics[#1]{../../output/regression/county_v2_ln_event_csdid/#2}}
\newcommand{\hospitalfig}[2][]{%
  \includegraphics[#1]{../../output/regression/hospital_v6/#2}}
\newcommand{\sym}[1]{\rlap{#1}}

\begin{document}

\begin{frame}{County-Level Empirical Specification}
\small
\begin{itemize}
  \item \textbf{Unit and sample:} county-by-year panel, 2012--2019.
  \item \textbf{Outcome:} total number of pharmacies in county $c$ and year $t$.
  \item \textbf{Treatment:}
  \[
    Reform_{ct}=\mathbf{1}\{t\geq PolicyYear_c\}.
  \]
\end{itemize}
\vspace{-0.4em}
\[
Y_{ct}=\beta Reform_{ct}+\gamma'X_{ct}
       +\alpha_c+\lambda_{p(c)\times t}+\varepsilon_{ct}.
\]
\begin{itemize}
  \item $X_{ct}$ includes log population and log mean nighttime lights.
  \item $\alpha_c$: county fixed effects; $\lambda_{p(c)\times t}$: province-by-year fixed effects.
  \item Standard errors are clustered at the county level.
  \item CSDID uses not-yet-treated counties as controls; event time $-1$ is omitted.
\end{itemize}
\end{frame}

\begin{frame}{County Total Pharmacies: Main Estimates}
\centering
\scriptsize
\begin{tabular}{lccc}
\toprule
 & TWFE & CSDID & Dependent variable \\
\midrule
Reform effect
  & $8.380^{***}$ & $7.311^{**}$ & Pharmacy count \\
  & $(2.245)$     & $(2.945)$    & \\
\midrule
County fixed effects       & Yes & Cohort-time ATT & \\
Province-by-year effects   & Yes & ---             & \\
Controls                   & Yes & Yes             & \\
Observations               & 19,274 & ---           & \\
\bottomrule
\end{tabular}

\vspace{1.0em}
\begin{itemize}
  \item Both estimators indicate an increase in the total number of pharmacies after reform.
  \item TWFE absorbs county and province-by-year fixed effects.
  \item CSDID is the preferred staggered-adoption comparison because it avoids already-treated controls.
\end{itemize}

\vfill
\tiny Notes: Standard errors in parentheses. $^{***}p<0.01$, $^{**}p<0.05$, $^{*}p<0.10$.
\end{frame}

\begin{frame}{County Total Pharmacies: TWFE Event Study}
\centering
\countyfig[width=0.82\textwidth]{twfe_event_county_pharmacy_count.png}

\vspace{-0.3em}
\small
Pre-treatment coefficients are jointly insignificant ($p=0.669$).
The post-reform estimates are positive and increase over event time.
\end{frame}

\begin{frame}{County Total Pharmacies: CSDID Event Study}
\centering
\countyfig[width=0.82\textwidth]{csdid_event_county_pharmacy_count.png}

\vspace{-0.3em}
\small
The dynamic CSDID profile supports a post-reform increase in county pharmacy supply.
Event time $-1$ is the reference period.
\end{frame}

\begin{frame}{County IV-DID: Specification and Estimates}
\small
\[
D_{ct}=Reform_{ct}\times
       \ln(PublicHospitals_{c,\mathrm{baseline}}).
\]
\[
D_{ct}\ \text{is instrumented by}\quad
Reform_{ct}\times Z_c,
\]
where $Z_c$ is either the 1959 epidemic stock or exposure in station-years during
the 1950s. The specification uses the same controls, county fixed effects,
province-by-year fixed effects, and county-clustered standard errors.

\vspace{0.5em}
\centering
\scriptsize
\begin{tabular}{lcc}
\toprule
 & Epidemic stock IV & Station-years IV \\
\midrule
IV-DID coefficient & $60.328^{***}$ & $68.753^{***}$ \\
                   & $(13.692)$     & $(15.549)$ \\
Kleibergen--Paap first-stage $F$ & 89.719 & 76.020 \\
\bottomrule
\end{tabular}

\vspace{0.8em}
\small
The instruments are strong in the reported first stage. The coefficient is the
effect of the reform intensity measure $D_{ct}$, not a one-unit binary DID effect.
\end{frame}

\begin{frame}{County IV-DID: Reduced-Form Event Studies}
\begin{columns}[T,onlytextwidth]
\column{0.49\textwidth}
\centering
\countyfig[width=\linewidth]{iv_rf_event_stock_count.png}

\scriptsize (a) 1959 epidemic stock

\column{0.49\textwidth}
\centering
\countyfig[width=\linewidth]{iv_rf_event_years_count.png}

\scriptsize (b) 1950s station-years exposure
\end{columns}

\vspace{0.4em}
\small
These dynamic reduced-form profiles diagnose the timing of the instrument--outcome
relationship. They should be interpreted cautiously where pre-reform coefficients
depart from zero.
\end{frame}

\begin{frame}{County Location Mechanism: Definition and Specification}
\small
\begin{itemize}
  \item Each geocoded pharmacy is assigned \textbf{once} to its nearest active public
        hospital within the same county:
  \[
    d_{jct}=\min_{h\in\mathcal{H}_{ct}}
            Distance(Pharmacy_{jct},Hospital_h).
  \]
  \item County outcomes are the number of uniquely assigned pharmacies with
        $d_{jct}\leq 2$ km and the number with $d_{jct}>2$ km.
  \item The 2 km split is the main mechanism measure; 500 m and 1 km thresholds
        are supporting robustness checks.
  \item The same county-level TWFE and CSDID specifications are re-estimated with
        these location outcomes.
\end{itemize}

\vspace{0.4em}
\[
N^{\leq 2km}_{ct}+N^{>2km}_{ct}
=N^{geocoded}_{ct}.
\]
\centering
\small
The identity holds in all 19,590 county-year records, with zero violations.
\end{frame}

\begin{frame}{County Location Mechanism: Static Estimates}
\centering
\scriptsize
\begin{tabular}{lcc}
\toprule
Outcome & TWFE & CSDID \\
\midrule
Nearest public hospital $\leq 2$ km
  & $12.498^{***}$ & $25.253^{***}$ \\
  & $(1.425)$      & $(0.837)$ \\
Nearest public hospital $>2$ km
  & $-4.118^{**}$  & $-17.942^{***}$ \\
  & $(2.053)$      & $(2.947)$ \\
\midrule
Sum of the two estimates
  & $8.380$        & $7.311$ \\
\midrule
Nearest public hospital $\leq 500$ m
  & $3.280^{***}$  & --- \\
Nearest public hospital $\leq 1$ km
  & $7.577^{***}$  & --- \\
\bottomrule
\end{tabular}

\vspace{0.9em}
\begin{itemize}
  \item The accounting identity reproduces the county total-pharmacy estimate.
  \item The signs suggest relocation or concentration toward public hospitals.
  \item Because the mechanism outcomes show pre-trends, these estimates are
        descriptive evidence rather than clean causal effects.
\end{itemize}

\vfill
\tiny Notes: County controls and fixed effects follow the main specification.
Standard errors are shown only for the 2 km main split.
\end{frame}

\begin{frame}{County Location Mechanism: TWFE Event Studies}
\begin{columns}[T,onlytextwidth]
\column{0.49\textwidth}
\centering
\countyfig[width=\linewidth]{twfe_event_mechanism_le2km.png}

\scriptsize (a) Nearest public hospital $\leq 2$ km

\column{0.49\textwidth}
\centering
\countyfig[width=\linewidth]{twfe_event_mechanism_gt2km.png}

\scriptsize (b) Nearest public hospital $>2$ km
\end{columns}

\vspace{0.4em}
\small
Joint pre-trend tests reject for both outcomes ($p<0.001$).
The figures are therefore informative about spatial patterns, not sufficient for
a causal mechanism claim.
\end{frame}

\begin{frame}{County Location Mechanism: CSDID Event Studies}
\begin{columns}[T,onlytextwidth]
\column{0.49\textwidth}
\centering
\countyfig[width=\linewidth]{csdid_event_mechanism_le2km.png}

\scriptsize (a) Nearest public hospital $\leq 2$ km

\column{0.49\textwidth}
\centering
\countyfig[width=\linewidth]{csdid_event_mechanism_gt2km.png}

\scriptsize (b) Nearest public hospital $>2$ km
\end{columns}

\vspace{0.4em}
\small
CSDID also displays pre-reform divergence. The location split is retained as
supplementary descriptive evidence.
\end{frame}

\begin{frame}{Hospital-Level Empirical Specification}
\small
\begin{itemize}
  \item \textbf{Unit and sample:} hospital-by-year panel, 2012--2019.
  \item \textbf{Outcome:} number of pharmacies within radius $r$ of hospital $i$;
        $r=2$ km is primary, while 1 km and 500 m are robustness checks.
  \item \textbf{Treatment contrast:} public versus private hospitals before and
        after the county reform.
\end{itemize}
\vspace{-0.4em}
\[
Y_{ict}(r)=\beta[Public_i\times Reform_{ct}]
 +\theta Reform_{ct}+\gamma NTL_{ict}(r)
 +\alpha_i+\lambda_{p(c)\times t}+\varepsilon_{ict}.
\]
\begin{itemize}
  \item $\alpha_i$: hospital fixed effects; $\lambda_{p(c)\times t}$:
        province-by-year fixed effects.
  \item Standard errors are clustered at the county level.
  \item Buffers may overlap: one pharmacy can be counted around multiple hospitals.
        This differs from the unique nearest-hospital county measure.
\end{itemize}
\end{frame}

\begin{frame}{Hospital Buffers: TWFE Estimates}
\centering
\scriptsize
\begin{tabular}{lccc}
\toprule
 & 2 km buffer & 1 km buffer & 500 m buffer \\
\midrule
Public hospital $\times$ reform
  & $-3.727^{***}$ & $-1.184^{***}$ & $-0.233^{***}$ \\
  & $(0.583)$      & $(0.211)$      & $(0.076)$ \\
\midrule
Hospital fixed effects       & Yes & Yes & Yes \\
Province-by-year effects     & Yes & Yes & Yes \\
Buffer nighttime-light control & Yes & Yes & Yes \\
Observations                 & 138,567 & 138,567 & 138,567 \\
\bottomrule
\end{tabular}

\vspace{1.0em}
\begin{itemize}
  \item Relative to private hospitals in the same reform environment, pharmacy
        counts fall more around public hospitals.
  \item The absolute magnitude declines monotonically as the buffer narrows.
  \item Event-study diagnostics are necessary because the static estimates alone
        do not establish parallel trends.
\end{itemize}

\vfill
\tiny Notes: Standard errors in parentheses. $^{***}p<0.01$.
\end{frame}

\begin{frame}{Hospital Buffers: TWFE Event Studies}
\begin{columns}[T,onlytextwidth]
\column{0.333\textwidth}
\centering
\hospitalfig[width=\linewidth]{twfe_event_pharm_count_2km.png}

\scriptsize (a) 2 km

\column{0.333\textwidth}
\centering
\hospitalfig[width=\linewidth]{twfe_event_pharm_count_1km.png}

\scriptsize (b) 1 km

\column{0.333\textwidth}
\centering
\hospitalfig[width=\linewidth]{twfe_event_pharm_count_500m.png}

\scriptsize (c) 500 m
\end{columns}

\vspace{0.5em}
\small
Joint pre-trend tests reject at 2 km ($p=0.014$), 1 km ($p=0.002$), and
500 m ($p=0.005$). The negative TWFE estimates should therefore be interpreted
with caution.
\end{frame}

\begin{frame}{Hospital Buffers: CSDID Estimates}
\centering
\scriptsize
\begin{tabular}{lccc}
\toprule
 & 2 km buffer & 1 km buffer & 500 m buffer \\
\midrule
Overall ATT
  & $-0.248$ & $-0.084$ & $0.072$ \\
  & $(0.343)$ & $(0.154)$ & $(0.063)$ \\
\midrule
Statistically significant at 5\% & No & No & No \\
Control group & \multicolumn{3}{c}{Not-yet-treated hospitals} \\
\bottomrule
\end{tabular}

\vspace{1.0em}
\begin{itemize}
  \item The CSDID point estimates are small and statistically insignificant at
        all three radii.
  \item The contrast with TWFE indicates sensitivity to staggered-treatment
        comparisons and pre-reform dynamics.
  \item The event studies provide the relevant diagnostic for interpreting these
        average treatment effects.
\end{itemize}
\end{frame}

\begin{frame}{Hospital Buffers: CSDID Event Studies}
\begin{columns}[T,onlytextwidth]
\column{0.333\textwidth}
\centering
\hospitalfig[width=\linewidth]{csdid_event_pharm_count_2km.png}

\scriptsize (a) 2 km

\column{0.333\textwidth}
\centering
\hospitalfig[width=\linewidth]{csdid_event_pharm_count_1km.png}

\scriptsize (b) 1 km

\column{0.333\textwidth}
\centering
\hospitalfig[width=\linewidth]{csdid_event_pharm_count_500m.png}

\scriptsize (c) 500 m
\end{columns}

\vspace{0.5em}
\small
The pre-reform dynamics, especially at event time $-2$, remain unstable.
Accordingly, the hospital-buffer estimates do not provide robust causal evidence
of a post-reform change.
\end{frame}

\end{document}
